% SPDX-License-Identifier: CC-BY-SA-4.0
% Author: Matthieu Perrin

\begingroup

\begin{exercice}[Mini-défi -- Le jeu de Fizz-buzz]\label{exo:lexing/automata/challenge}
  
  Dans le jeu Fizz-buzz, les joueurs comptent à voix haute à tour de rôle, en remplaçant
  tout nombre divisible par trois par le mot \og fizz \fg,
  tout nombre divisible par cinq par le mot \og buzz \fg,
  et tout nombre divisible par quinze par le mot \og fizz-buzz \fg.
  La séquence commence donc par
  \begin{center}
    \og 1, 2, fizz, 4, buzz, fizz, 7, 8, fizz, buzz, 11, fizz, 13, 14, fizz-buzz, etc \fg
  \end{center}
  
  On travaille sur l'alphabet $\{0,1\}$.
  Le but de cet exercice est de concevoir un automate fini capable de catégoriser
  les nombres selon leur divisibilité par $3$, $5$ et $15$. 
  
  Un mot $w$ (sans zéro initial inutile) code en binaire un entier $N(w)$.
  On veut concevoir un AFD qui, à la fin de la lecture,
  \begin{itemize}
  \item s'arrête dans un état de \textbf{catégorie} \emph{fizz-buzz} si $15 \mid N(w)$,
  \item s'arrête dans un état de \textbf{catégorie} \emph{fizz} si $3 \mid N(w)$ et $15 \nmid N(w)$,
  \item s'arrête dans un état de \textbf{catégorie} \emph{buzz} si $5 \mid N(w)$ et $15 \nmid N(w)$,
  \item sinon s'arrête dans un état non accepteur (\emph{autre}).
  \end{itemize}
  
  \begin{question}
  \item Proposer un AFD satisfaisant ces exigences (décrire $Q,\ \Sigma,\ \delta,\ q_0$, et partitionner $F$ en trois \emph{catégories d’acceptation} : \textsc{fizz}, \textsc{buzz}, \textsc{fizz-buzz}).
  \item (Optionnel) Donner un petit tableau de transitions représentatif et/ou un schéma partiel.
  \end{question}
  
\end{exercice}

\begin{correction}
  On représente l'automate comme un tableau
  
  \begin{tabular}{|c|c|c|c|}
    \hline
    état & catégorie & 0 & 1 \\
    \hline
    0  & fizz-buzz & 0  & 1  \\ 
    1  &           & 2  & 3  \\ 
    2  &           & 4  & 5  \\ 
    3  & fizz      & 6  & 7  \\ 
    4  &           & 8  & 9  \\ 
    5  & buzz      & 10 & 11 \\ 
    6  & fizz      & 12 & 13 \\ 
    7  &           & 14 & 0  \\ 
    8  &           & 1  & 2  \\ 
    9  & fizz      & 3  & 4  \\ 
    10 & buzz      & 5  & 6  \\ 
    11 &           & 7  & 8  \\ 
    12 & fizz      & 9  & 10 \\ 
    13 &           & 11 & 12 \\ 
    14 &           & 13 & 14 \\ 
    \hline
  \end{tabular}
\end{correction}

\endgroup
\endinput
